% =====================================================================
% PASTE FILE -- Approach / Results / What Did Not Work / Conclusions,
% adaptate din report/report.tex la formatul de workshop (MRL 2026).
% Se copiaza dupa Related Work in acl_latex.tex, inainte de Limitations.
%
% Fata de report/report.tex:
%   - cheile de citare noi (vezi maparea din paste_bib_additions.bib);
%   - metrica definita explicit si marcata ca metrica a competitiei;
%   - sectiune noua 4.5 "What was tuned" -- divulgarea completa;
%   - subsectiune noua 5.2 -- evaluarea pe cele 14 exemple RO;
%   - tab:fail rescris pe doua blocuri: cele doua metrici in-source nu mai
%     sunt amestecate (0.515 era EN->ML, 0.586 era transfer_avg);
%   - "87% din scor" scos: nu se reconcilia cu propriul tabel (e 40%);
%   - fara link de GitHub, fara TODO de author contributions.
%
% NECESITA IN PREAMBUL: \usepackage{booktabs} si \graphicspath{{figures/}}
% -- vezi paste_preamble.tex. (enumitem NU e folosit nicaieri.)
% =====================================================================

\section{Approach}
\label{section:approach}

\subsection{Task, data and metric}

Given a target word in a Romanian sentence, the task is to predict a
complexity score in $[0,1]$, where $0$ is a fully familiar word and $1$
a maximally difficult one. The test set contains $1587$ target words.
Fourteen Romanian items carry gold complexity labels; they are too few
to train or validate on, and we use them only as an independent check
(Section~\ref{sec:14items}). Aggregated eye-tracking measurements over
the same texts are distributed with the task, so gaze is available for
the target words without any annotation effort of our own.

Systems are ranked by
\begin{equation}
  \label{eq:metric}
  \mathrm{score} = 50\,\bigl(|r| + \max(0, R^2)\bigr),
\end{equation}
where $r$ is the Pearson correlation with the gold scores and $R^2$ the
coefficient of determination. This is the metric defined by the shared
task, not a standard LCP measure, and it conflates two different
things: $|r|$ rewards ranking the words correctly, while $R^2$ rewards
putting them on the right absolute scale. We report it because the
leaderboard returns only this composite --- the gold labels for the
test set are not released, so we cannot decompose our own submissions
into $r$ and $R^2$. Wherever labels are available we report $r$
separately, and we take $r$ as the primary measure for the reason given
in Section~\ref{sec:tuning}.

\subsection{Setup and in-source validation}

The system uses XLM-RoBERTa-base \citep{conneau2020xlmr} and RoLlama3-8B
\citep{masala-etal-2024-vorbesti} in 4-bit quantisation, with
\texttt{transformers}, \texttt{scikit-learn} and \texttt{wordfreq}.

With no Romanian labels to validate on, design decisions were made
\emph{in-source}: a regressor is trained on one set of source languages
and evaluated on another, with cross-validation inside each source. We
report two such quantities, and they are not interchangeable.
$r_{\mathrm{EN}\rightarrow\mathrm{RO}^{*}}$ is the correlation obtained
when training on English CompLex \citep{shardlow-etal-2021-semeval} and
testing on the Romance portion of MLSP \citep{shardlow-etal-2024-bea};
$\bar{r}_{\mathrm{transfer}}$ averages that with the reverse direction.
Built on thousands of real labels, either check tells us whether a
configuration is sound before it is applied to Romanian. Its limits are
the subject of Section~\ref{sec:negative}.

\subsection{The four signals}

Each signal estimates complexity in a different way, and none of them
requires a Romanian complexity annotation.

\paragraph{(A) Cross-lingual transfer.}
Complexity is learned from English and Romance annotations
\citep{shardlow-etal-2021-semeval, shardlow-etal-2024-bea} and carried
over to Romanian through the shared representation space of
XLM-RoBERTa, which places translation-equivalent words close together
(Section~\ref{sec:related}). A Ridge and HistGradientBoosting blend
reads each word's embedding --- the mean of the last four layers,
centred per language --- together with a few surface cues: frequency,
length, syllable count and capitalisation.

One Romanian-specific correction accounts for the largest single gain
in the system. Because Romanian is heavily inflected, many word forms
are absent from the frequency list and would be scored as maximally
rare: of the $1587$ target tokens, $11.7\%$ have no entry at all, and
for $36.0\%$ the lemma is more frequent than the inflected form, by
$0.86$ Zipf units on average. We therefore use
$\max\bigl(\mathrm{zipf}(\mathrm{form}), \mathrm{zipf}(\mathrm{lemma})\bigr)$
in place of the form frequency. Section~\ref{sec:14items} identifies the
condition under which this substitution is the wrong thing to do.

\paragraph{(B) Language-model rating.}
RoLlama3 is asked to rate the target word in context on a $1$--$5$
scale. We read the answer from the probabilities the model assigns to
the five digit tokens rather than from the text it generates, which
turns a discrete choice into a smooth expected value.

\paragraph{(C) Surprisal.}
Surprisal, $-\log P(w \mid \mathrm{context})$, measures how unexpected
the word is; under surprisal theory an unexpected word is harder to
process \citep{hale2001probabilistic}. It comes from the same forward
pass as (B) but is a raw internal statistic, which is why it is nearly
independent of the rating ($r = 0.07$).

\paragraph{(D) Eye-tracking.}
Harder words are fixated longer and re-read more often
\citep{rayner1998eye}. We compress the eight per-word effort measures
into a single reading-effort score with PCA. The first component
explains $62\%$ of the variance and every measure loads on it with the
same sign (Table~\ref{tab:pca}), so the component is interpretable as a
single ``more effort $=$ more complex'' axis.

\subsection{Fusing the signals}

The four signals live on different scales, so each is standardised and
they are combined as a weighted average. Signal A is itself an ensemble
of the English-source and Romance-source transfers, averaged after
standardisation. The fused score is then mapped onto the scale of the
test distribution,
\begin{equation}
  \label{eq:calib}
  \hat{y} = \mathrm{clip}\bigl(\mu + \sigma \cdot z(s),\; 0,\; 1\bigr),
\end{equation}
with $\mu = 0.27$ and $\sigma = 0.125$. Per-signal weights are in
Table~\ref{tab:signals}.

\subsection{What was tuned, and what it costs}
\label{sec:tuning}

A zero-label setting has no development set, and selecting anything
against the evaluation score is a known way to inflate it
\citep{gorman-bedrick-2019-need}. We state exactly what this affects.

Five numbers in the system were chosen against the public leaderboard
score: the three fusion weights and the two constants of
Equation~\ref{eq:calib}. The weights were searched on a coarse grid;
for the eye-tracking weight the full sweep was
$w_D \in \{0, 0.10, 0.20, 0.30\} \rightarrow 68.90, 69.66, 69.74,
69.00$. Everything else --- the encoder, the layer, the feature set,
the regressor, the lemma correction --- was decided in-source, without
reference to the leaderboard.

Two things bound the damage. First, the calibration constants cannot
affect $r$ at all. Equation~\ref{eq:calib} is affine up to the clip,
and the clip is inactive on our predictions, whose range is
$[0.014, 0.845]$; Pearson $r$ is invariant under affine rescaling, so
$\mu$ and $\sigma$ --- the two quantities that most directly encode
information about the test distribution --- move only the $R^2$ term of
Equation~\ref{eq:metric}. This is why we treat $r$ as the primary
measure. Second, the response surface is flat: the eye-tracking sweep
spans $0.84$ points end to end, and on the 14 Romanian items the tuned
weights are worth $0.019$ of Pearson $r$ over uniform weights
(Section~\ref{sec:14items}). The one number that must be read as an
upper bound is the $+0.84$ attributed to eye-tracking in
Table~\ref{tab:ablation}: it is the maximum of the sweep that selected
$w_D$, and no independent estimate of it is available.

\section{Results}
\label{sec:results}

\subsection{How much each signal contributes}

Table~\ref{tab:ablation} builds the system one signal at a time.
Restricted to what the shared task itself provides --- surface features
and eye-tracking --- the system reaches $41.76$; the imported signals
take it to $69.74$, so roughly $40\%$ of the final score is
attributable to knowledge brought in from outside Romanian. Surface
features alone reach $10.27$. The single largest step is adding the
cross-lingual and language-model signals together ($+22.76$); the
lemma-frequency correction is the largest single gain after that
($+4.09$), and surprisal and eye-tracking each add a small amount.

\subsection{Evaluation on the 14 Romanian items}
\label{sec:14items}

The fourteen labelled Romanian items are the only real Romanian
supervision in existence for this task. They were never used for
training and took no part in weight selection, which was carried out
against the leaderboard; they were consulted during development as a
sanity check, so they are not a pristine test set either, but they are
independent of the tuning audited in Section~\ref{sec:tuning}.
Table~\ref{tab:ro14} reports Pearson $r$ on them. Two caveats apply
throughout: $n = 14$, so the confidence intervals are wide; and only
signals A and B are available on these items, because surprisal was
never computed for them and only $2$ of the $14$ carry eye-tracking.

The fused signal reaches $r = 0.913$, and a permutation test rejects
chance at $p < 0.0001$ despite the sample size. Uniform weights give
$0.894$: the tuning is worth $0.019$, which is the estimate quoted in
Section~\ref{sec:tuning}. With the calibration constants of
Equation~\ref{eq:calib} the composite metric on these items is $73.53$
($R^2 = 0.558$), close to the $69.74$ obtained on the full test set and
in the same direction, which is weak evidence that the leaderboard
figure is not inflated.

\paragraph{A scope condition for the lemma correction.}
The first two rows of Table~\ref{tab:ro14} are the frequency baselines,
and they invert. Raw form frequency correlates $-0.606$ with
complexity, the expected direction; after the lemma correction the
correlation is $+0.092$. The reversal comes from a single item:
\emph{putu}, a simple-perfect verb form, is the most complex word in the
set ($0.545$) and the correction raises its frequency estimate by
$2.76$ Zipf units, turning the hardest word into one of the easiest.
Excluding it, the two baselines are $-0.465$ and $-0.476$, so on the
remaining thirteen items the correction is neutral to mildly helpful.
The correction assumes that an inflected form inherits its difficulty
from its lemma. That assumption fails exactly when the inflection is
itself the source of difficulty --- archaic, literary or otherwise
marked forms --- and Romanian simple-perfect verbs are such a case. The
$+4.09$ measured over $1587$ items says the assumption usually holds;
this item says it is an assumption, not an identity.

\subsection{What drives the prediction}

A SHAP analysis \citep{lundberg2017shap} over named features
(Figure~\ref{fig:beeswarm}) places
frequency first, the language-model rating second and surprisal third;
the surface features barely matter. Table~\ref{tab:signals} gives the
same ordering at the signal level, and it agrees with the tuned
weights --- SHAP describes the fitted model, it does not choose the
weights. The signals are complementary rather than redundant
(Figure~\ref{fig:corr}): the rating and surprisal are nearly
independent, while the cross-lingual signal tracks frequency closely
enough to be read as a frequency detector with semantics attached.

\subsection{Eye-tracking contributes through its late measures}

Eye-tracking helps because it carries something the lexical signals do
not, and the division is the one predicted by the reading literature.
Partialling out frequency and length from each gaze measure over the
$1509$ test words that have gaze data, the early measures --- first
fixation duration, first-pass reading time, skip rate --- retain almost
no relation to the system's prediction ($\bar{r} = 0.010$) or to
surprisal ($\bar{r} = -0.027$): what they encode is lexical access,
which signal A already knows. The late measures --- re-reading time,
incoming regressions, refixation rate --- retain $\bar{r} = 0.243$ and
$\bar{r} = 0.177$ respectively. Late gaze is where the
contextual-integration difficulty lives, and that is the part of the
eye-tracking signal that is not already in the other three. As a
separate check on the transfer regressor itself, on held-out English
words it sorts complexity into easy, medium and hard bands with a
clearly diagonal confusion matrix (Figure~\ref{fig:confusion}).

\begin{figure}[tbp]
\centering
\includegraphics[width=\columnwidth]{beeswarm.png}
\caption{SHAP over named features (top 8 shown). Frequency dominates, then the LLM rating and surprisal. The
separate blue cluster on the frequency row is the group of very rare / out-of-vocabulary words, which push the
predicted complexity up --- the expected behavior. Red = high feature value.}
\label{fig:beeswarm}
\end{figure}

\begin{figure}[tbp]
\centering
\includegraphics[width=\columnwidth]{fig_corr_matrix.png}
\caption{Correlation between the four signals and frequency on the Romanian test set. The rating (B) and surprisal
(C) are nearly independent; the cross-lingual signal (A) tracks frequency closely.}
\label{fig:corr}
\end{figure}

\begin{figure}[tbp]
\centering
\includegraphics[width=0.8\columnwidth]{fig_confusion.png}
\caption{In-source check: on held-out English words, the regressor sorts complexity into easy/medium/hard bands
with a clearly diagonal confusion matrix.}
\label{fig:confusion}
\end{figure}

\begin{table}[tbp]
\centering\small
\begin{tabular}{lcc}
\toprule
\textbf{Configuration} & \textbf{Score} & \textbf{$\Delta$} \\
\midrule
Predict a constant (the mean) & 0.00 & --- \\
Surface features only (zero-shot RO) & 10.27 & +10.27 \\
\;+ eye-tracking (no imported res.) & 41.76 & +31.49 \\
\;+ cross-lingual A + LLM B & 64.52 & +22.76 \\
\;+ surprisal C & 64.81 & +0.29 \\
\;+ lemma-frequency fix (in A) & 68.90 & +4.09 \\
\;+ eye-tracking PCA D & \textbf{69.74} & +0.84$^{\dagger}$ \\
\bottomrule
\end{tabular}
\caption{Shared-task score (Equation~\ref{eq:metric}), built one component at a time; $\Delta$ is the marginal
contribution of each component. Predicting one constant gives no correlation, hence $0$. The first two rows use
only resources the shared task provides, so the imported signals account for $(69.74-41.76)/69.74 \approx 40\%$ of
the final score. $^{\dagger}$This delta is the maximum of the sweep that selected $w_D$ and should be read as an
upper bound (Section~\ref{sec:tuning}).}
\label{tab:ablation}
\end{table}

\begin{table}[tbp]
\centering\small
\begin{tabular}{lc}
\toprule
\textbf{Predictor on the 14 items} & \textbf{Pearson $r$} \\
\midrule
Zipf frequency, inflected form & $-0.606$ \\
Zipf frequency, lemma-corrected & $+0.092$ \\
\midrule
A cross-lingual, English source & $0.891$ \\
A cross-lingual, Romance source & $0.759$ \\
A ensemble & $0.838$ \\
B language-model rating & $0.615$ \\
\midrule
A+B, uniform weights (untuned) & $0.894$ \\
A+B, tuned weights & $\mathbf{0.913}$ \\
\bottomrule
\end{tabular}
\caption{Pearson $r$ against the gold labels of the 14 Romanian items, which took no part in weight selection.
Complexity falls as frequency rises, so a negative $r$ is the expected direction for the frequency rows and a
positive one for the model signals. Bootstrap $95\%$ CIs over $10{,}000$ resamples: $[0.661, 0.967]$ for the
untuned fusion and $[0.663, 0.975]$ for the tuned one; a permutation test over $20{,}000$ shuffles gives
$p < 0.0001$ for the latter. Signals C and D are absent: surprisal was not computed for these items and only $2$
of the $14$ have eye-tracking.}
\label{tab:ro14}
\end{table}

\begin{table}[tbp]
\centering\small
\begin{tabular}{lcc}
\toprule
\textbf{Signal} & \textbf{SHAP importance} & \textbf{weight} \\
\midrule
A cross-lingual & 0.061 & 44\% \\
B LLM rating & 0.034 & 29\% \\
D eye-tracking & 0.017 & 20\% \\
C surprisal & 0.006 & 6\% \\
\bottomrule
\end{tabular}
\caption{Per-signal SHAP importance and the weight each signal gets in the final average. The two orders agree;
SHAP describes the model, it does not choose the weights.}
\label{tab:signals}
\end{table}

\begin{table}[tbp]
\centering\small
\begin{tabular}{lc}
\toprule
\textbf{Eye-tracking metric} & \textbf{PC1 loading} \\
\midrule
Total fixation time (TFT) & 0.44 \\
Total fixation count (TFC) & 0.43 \\
First-pass reading time (FPRT) & 0.42 \\
First-pass fixation count (FPFC) & 0.39 \\
Re-reading time (RRT) & 0.35 \\
First-fixation duration (FFD) & 0.29 \\
Regression-path duration (RPD\_inc) & 0.28 \\
Incoming regressions (TRC\_in) & 0.13 \\
\bottomrule
\end{tabular}
\caption{Weight of each eye-tracking metric in the reading-effort score (loadings on PC1, which explains $62\%$ of
the variance). All loadings are positive: one ``more effort = more complex'' axis.}
\label{tab:pca}
\end{table}

\section{What Did Not Work}
\label{sec:negative}

Table~\ref{tab:fail} lists the attempts that failed, and two of them
carry a lesson beyond this system.

\paragraph{Fine-tuning is worse than freezing.}
Every form of fine-tuning we tried --- full fine-tuning and LoRA, over
a sweep of learning rates, epochs and ranks --- transferred worse than
the frozen encoder ($0.461$ and $0.462$ against $0.515$). With no
target-language labels, fine-tuning can only adapt to the source
language, and in doing so it degrades the cross-lingual alignment that
the transfer depends on. This is the behaviour \citet{peters-etal-2019-tune}
report when the target task is distant from the source, and it is worth
stating plainly because the default move in a cross-lingual pipeline is
to fine-tune \citep{pires-etal-2019-multilingual}.

\paragraph{In-source validation does not predict transfer to Romanian.}
Three configurations beat the frozen XLM-R baseline on in-source
transfer and then lost on Romanian: a different encoder (e5), mid-layer
embeddings, and a Lasso regressor. The in-source check measures
transfer between English and Romance --- related, well-resourced
languages --- and that is evidently not the same problem as transfer
into Romanian. In a setting with no target-language development set
this is the uncomfortable conclusion: the validation signal available
is not the validation signal needed.

\paragraph{Prompting and extra features.}
Chain-of-thought and few-shot prompting added noise rather than signal;
at the $8$B scale chain-of-thought is not yet reliably helpful
\citep{wei2022cot}. Rarity measures and a range of part-of-speech,
named-entity and WordNet features proved redundant with frequency,
which is consistent with the SHAP picture in
Figure~\ref{fig:beeswarm}.

\begin{table}[tbp]
\centering\small
\begin{tabular}{lcc}
\toprule
\textbf{Attempt} & \textbf{in-source} & \textbf{score} \\
\midrule
\multicolumn{3}{l}{\emph{Fine-tuning} ($r_{\mathrm{EN}\rightarrow\mathrm{RO}^{*}}$)} \\
Frozen XLM-R (this work) & 0.515 & --- \\
Full fine-tune (best of sweep) & 0.461 & --- \\
LoRA (best of sweep) & 0.462 & --- \\
\midrule
\multicolumn{3}{l}{\emph{Representation and regressor} ($\bar{r}_{\mathrm{transfer}}$)} \\
Frozen XLM-R (this work) & 0.551 & 69.74 \\
Encoder swap (e5) & 0.586 & 63.00 \\
Mid-layer embeddings & 0.556 & 64.49 \\
Lasso regressor & 0.569 & 64.28 \\
\midrule
\multicolumn{3}{l}{\emph{Prompting}} \\
LLM few-shot prompt & --- & 63.31 \\
\bottomrule
\end{tabular}
\caption{Failed attempts. The two blocks use different in-source measures and are not comparable across blocks:
$r_{\mathrm{EN}\rightarrow\mathrm{RO}^{*}}$ is the one-directional English-to-Romance transfer correlation, and
$\bar{r}_{\mathrm{transfer}}$ averages both directions. Within the second block, every variant that beats the
frozen baseline in-source loses on the shared-task score.}
\label{tab:fail}
\end{table}

\section{Conclusions}
\label{section:conclusions}

The absence of Romanian training data did not merely constrain the
design, it determined it: with no labels to fit and no development set
to select on, the only remaining move is to import complexity knowledge
from elsewhere and combine several independent imports. Four such
signals --- a frozen multilingual encoder, a language-model rating,
surprisal, and gaze --- reach $69.74$ on the shared-task metric and
$r = 0.913$ on the labelled Romanian items released with the task,
without a single target-language annotation.

Three findings generalise beyond Romanian. A lemma-frequency correction
is worth more than any other single component ($+4.09$) for a
morphologically rich language, subject to the scope condition of
Section~\ref{sec:14items}. Fine-tuning a multilingual encoder without
target labels is actively harmful, and freezing it is the stronger
default. And in-source validation, the only validation a zero-label
setting affords, does not reliably predict transfer to the target
language --- which is the finding we would most want future work in
this setting to take seriously.
